% ============================================================
%  A Systematic Benchmark of Persistent Homology Pipelines
%  for Classification
% ============================================================

\documentclass[12pt, a4paper]{report}

\usepackage[T1]{fontenc}
\usepackage[utf8]{inputenc}

\usepackage[a4paper, top=2.5cm, bottom=2.5cm,
            left=3.0cm, right=2.5cm]{geometry}
\usepackage{setspace}
\onehalfspacing

\usepackage{amsmath}
\usepackage{amssymb}
\usepackage{amsthm}
\usepackage{mathtools}

\newcommand{\VR}{\mathrm{VR}}
\newcommand{\Dgm}{\mathrm{Dgm}}
\newcommand{\birth}{\mathrm{b}}
\newcommand{\death}{\mathrm{d}}
\newcommand{\pers}{\mathrm{p}}
\newcommand{\Betti}{\beta}
\newcommand{\R}{\mathbb{R}}
\newcommand{\Z}{\mathbb{Z}}
\newcommand{\N}{\mathbb{N}}
\newcommand{\E}{\mathbb{E}}
\DeclareMathOperator{\im}{im}
\DeclareMathOperator{\coker}{coker}

\usepackage{graphicx}
\usepackage{float}
\usepackage{caption}
\captionsetup{font=small, labelfont=bf}
\graphicspath{{./figures/}}

\usepackage{booktabs}
\usepackage{tabularx}
\usepackage{array}

\usepackage{listings}
\usepackage{xcolor}

\definecolor{codegray}{rgb}{0.5,0.5,0.5}
\definecolor{codeblue}{rgb}{0.15,0.35,0.65}
\definecolor{codegreen}{rgb}{0.1,0.5,0.1}
\definecolor{codebg}{rgb}{0.97,0.97,0.97}

\lstdefinestyle{python}{
  language        = Python,
  backgroundcolor = \color{codebg},
  basicstyle      = \ttfamily\footnotesize,
  keywordstyle    = \color{codeblue}\bfseries,
  commentstyle    = \color{codegray}\itshape,
  stringstyle     = \color{codegreen},
  numberstyle     = \tiny\color{codegray},
  numbers         = left,
  stepnumber      = 1,
  breaklines      = true,
  frame           = single,
  captionpos      = b
}

\lstdefinestyle{yaml}{
  language        = {},
  backgroundcolor = \color{codebg},
  basicstyle      = \ttfamily\footnotesize,
  keywordstyle    = \color{codeblue}\bfseries,
  commentstyle    = \color{codegray}\itshape,
  stringstyle     = \color{codegreen},
  numberstyle     = \tiny\color{codegray},
  numbers         = left,
  breaklines      = true,
  frame           = single,
  captionpos      = b
}

\usepackage[colorlinks=true,
            linkcolor=blue!60!black,
            citecolor=blue!60!black,
            urlcolor=blue!60!black]{hyperref}

\theoremstyle{definition}
\newtheorem{definition}{Definition}[chapter]
\newtheorem{example}{Example}[chapter]

\theoremstyle{plain}
\newtheorem{theorem}{Theorem}[chapter]
\newtheorem{lemma}[theorem]{Lemma}
\newtheorem{corollary}[theorem]{Corollary}

\theoremstyle{remark}
\newtheorem{remark}{Remark}[chapter]


\begin{document}

% ════════════════════════════════════════════════════════════
%  TITLE PAGE
% ════════════════════════════════════════════════════════════
\begin{titlepage}
  \centering
  \vspace*{3cm}
  {\LARGE\bfseries A Systematic Benchmark of Persistent\par}
  \vspace{0.2cm}
  {\LARGE\bfseries Homology Pipelines for Classification\par}
  \vspace{2.5cm}
  {\large Zachariah Attila Markusson\par}
  \vspace{2cm}
  {\normalsize August 2026\par}
  \vfill
  {\small Independent Research\par}
\end{titlepage}

% ════════════════════════════════════════════════════════════
%  ABSTRACT
% ════════════════════════════════════════════════════════════
\begin{abstract}

Persistent homology classification pipelines have three stages ---
filtration, vectorization, and classifier --- each with multiple
available methods. Which stage matters most? The literature's most
widely cited answer, from Conti et al.\ \cite{conti2022}, is that filtration
choice dominates: switching filtrations on MNIST images swung
accuracy from 18\% to 94\%. We show that this finding does not
generalise. On ECG200 time series with Takens embedding,
vectorization choice dominates classification accuracy with a
marginal range of 6.1 percentage points (across 7 vectorizers,
95\% CI [4.8, 7.4]), while filtration choice contributes only
0.7pp (across 4 filtrations, 95\% CI [0.2, 1.1]).
Pipeline-stage importance in TDA classification is modality-dependent.
The sweep comprised 4 filtrations, 7 vectorizers, and 4 classifiers
(the framework is extensible to 7 filtrations, 11 vectorizers, and 4
classifiers via YAML-only changes).

We demonstrate this on one synthetic (sphere/torus at four noise
levels) and two real-world datasets (ECG200, MNIST binary), with
616 completed configurations (of 672 possible; 56 excluded where
the point-cloud filtrations Weak Alpha and Sparse Rips are
incompatible with image data). The topological signal on
sphere/torus point clouds survives Gaussian noise up to $\sigma = 0.30$
at 99.85\% mean accuracy across the 112 configurations at that level:
classification accuracy survives noise levels where the stability
bound's worst-case guarantee would permit degradation, and the
$\Betti_1 = 0$ vs.\ $\Betti_1 = 2$ distinction remains classifiable.
On MNIST, cubical filtration outperforms Vietoris-Rips
by 1.8 percentage points (98.0\% vs.\ 96.25\%), confirming that
filtration choice matters for image data. The Alpha complex is
11--21\% faster than Vietoris-Rips at identical accuracy on
100-point 3D clouds under the giotto-tda implementation. Simple
persistence statistics match kernel-based vectorizers on real data
despite requiring zero hyperparameters. All code, configuration, and
result schemas are provided; 616 completed configurations across
6 dataset instances, 4 filtrations, 7 vectorizers, and 4 classifiers.

\end{abstract}

\tableofcontents
\listoffigures
\listoftables
\newpage

% ════════════════════════════════════════════════════════════
%  CHAPTER 1 -- INTRODUCTION AND MOTIVATION
% ════════════════════════════════════════════════════════════
\chapter{Introduction and Motivation}
\label{chap:intro}

% ────────────────────────────────────────────────────────────
\section{Topological Data Analysis and the Pipeline Problem}
% ────────────────────────────────────────────────────────────

Topological Data Analysis (TDA) extracts shape invariants from data:
connected components, cycles, and voids that persist across multiple
scales. Unlike traditional machine learning which operates on raw
coordinates or pixel intensities, TDA captures the underlying
``shape'' of the data --- information invariant under stretching and
bending. Persistent homology, the workhorse of TDA, tracks the birth
and death of topological features as a scale parameter varies.

Practitioners face a combinatorial problem: the TDA classification
pipeline has three stages --- filtration, vectorization, and
classification --- each with 5--15 available methods (surveys:
Hatwar and Thangaraj \cite{hatwar2026}). The literature's
most widely cited answer to ``which stage matters most?'' comes from
Conti et al.\ \cite{conti2022}, who showed that filtration choice on MNIST
images can swing accuracy from 18\% to 94\% --- a finding that has
been interpreted as a general principle: ``filtration dominates.''

This paper argues that principle does not generalise. Pipeline-stage
importance is modality-dependent. We demonstrate this using a modular
benchmarking framework that varies all three pipeline stages
simultaneously in a controlled factorial sweep: 6 dataset instances
$\times$ 4 filtrations (Vietoris-Rips, Alpha, Sparse Rips, Cubical)
$\times$ 7 vectorizers $\times$ 4 classifiers --- 672 possible, 616
completed (56 excluded where point-cloud filtrations are incompatible
with image data). The framework is extensible to 7 filtrations, 11
vectorizers, and 4 classifiers via YAML-only configuration changes;
standardised harnesses of this kind are the direction the community
has moved towards \cite{telyatnikov2024}.

% ────────────────────────────────────────────────────────────
\section{The Persistent Homology Pipeline}
% ────────────────────────────────────────────────────────────

Whether persistent homology features help classification at all has
itself been questioned \cite{turkes2022}; this benchmark assumes the
features are informative and asks which pipeline stage extracts them
most effectively.

A persistent homology classification pipeline transforms raw data
through three stages:

\begin{enumerate}
  \item \textbf{Filtration.} Raw data is converted into a nested
    sequence of simplicial complexes parameterised by a scale
    $\varepsilon$. Topological features appear (birth at $\birth$)
    and disappear (death at $\death$). Output: a \emph{persistence
    diagram}, a multiset of $(\birth, \death)$ pairs.
  \item \textbf{Vectorization.} Diagrams are multisets of variable
    cardinality and cannot be fed directly into classifiers.
    Vectorization methods map each diagram to a fixed-length feature
    vector.
  \item \textbf{Classification.} The vectorized features are passed
    to a standard classifier which learns a decision boundary.
\end{enumerate}

\begin{figure}[H]
  \centering
  % [IMAGE: Horizontal pipeline diagram.]
  \caption{The three-stage persistent homology classification pipeline.}
  \label{fig:pipeline_diagram}
\end{figure}

% ────────────────────────────────────────────────────────────
\section{Contributions}
% ────────────────────────────────────────────────────────────

This paper demonstrates that pipeline-stage importance in TDA
classification is modality-dependent, using a modular benchmarking
framework built for that purpose. Specifically:

\begin{enumerate}
  \item \textbf{Finding.} On ECG200 time series with Takens embedding,
    vectorization dominates classification accuracy (6.1pp marginal
    range across 7 vectorizers, 95\% CI [4.8, 7.4]) while filtration
    is statistically neutral (0.7pp across 4 filtrations,
    95\% CI [0.2, 1.1]).
    This qualifies the widely cited claim that filtration always
    dominates --- that finding holds for image data but does not
    generalise. Pipeline-stage importance is fundamentally
    modality-dependent.
  \item \textbf{Framework.} The finding was produced by a benchmarking
    framework that varies all three pipeline stages simultaneously:
    4 filtrations, 7 vectorizers, and 4 classifiers executed in a
    616-configuration factorial sweep (of 672 possible; 56 exclusions
    where point-cloud filtrations are incompatible with image data),
    with YAML-driven configuration
    and normalised SQLite storage. The framework is extensible to
    7 filtrations, 11 vectorizers, and 4 classifiers via YAML-only
    changes. Full configuration and schema in Appendices~A--B; code
    implementations in Appendix~C.
\end{enumerate}

\noindent Chapter~2 develops the mathematical machinery needed to
understand why these pipeline stages interact as they do.

\newpage

% ════════════════════════════════════════════════════════════
%  CHAPTER 2 -- MATHEMATICAL FOUNDATIONS OF PERSISTENT TOPOLOGY
% ════════════════════════════════════════════════════════════
\chapter{Mathematical Foundations of Persistent Topology}
\label{chap:foundations}

This chapter develops the algebraic and geometric machinery
underlying persistent homology. We progress from simplicial complexes
through chain groups, homology vector spaces, filtrations, and
persistence modules, concluding with the stability theorem and
vectorization methods that convert diagrams into classifier-ready
feature vectors.

% ────────────────────────────────────────────────────────────
\section{Simplicial Complexes}
\label{sec:simplicial}
% ────────────────────────────────────────────────────────────

\begin{definition}[Simplex]
  A \emph{$k$-simplex} $\sigma = [v_0, v_1, \dots, v_k]$ is the
  convex hull of $k+1$ affinely independent points in $\R^d$
  ($d \ge k$). A 0-simplex is a vertex, a 1-simplex an edge, a
  2-simplex a triangle, a 3-simplex a tetrahedron. A \emph{face}
  of $\sigma$ is the convex hull of any non-empty subset of its
  vertices.
\end{definition}

\begin{definition}[Simplicial Complex]
  A \emph{simplicial complex} $K$ is a finite collection of simplices
  such that:
  \begin{enumerate}
    \item If $\sigma \in K$ and $\tau$ is a face of $\sigma$, then
      $\tau \in K$ (closure under taking faces).
    \item If $\sigma, \tau \in K$, then $\sigma \cap \tau$ is either
      empty or a face of both.
  \end{enumerate}
\end{definition}

The \emph{$k$-skeleton} $K^{(k)}$ is the subcomplex of all simplices
of dimension at most $k$. The 1-skeleton is a graph in the usual
sense.

% ────────────────────────────────────────────────────────────
\section{Filtrations and Complex Constructions}
\label{sec:filtrations}
% ────────────────────────────────────────────────────────────

\begin{definition}[Filtration]
  A \emph{filtration} of a simplicial complex $K$ is a nested
  sequence of subcomplexes indexed by a real parameter
  $\varepsilon \ge 0$:
  \begin{equation}
    \emptyset = K_{\varepsilon_0} \subseteq K_{\varepsilon_1}
    \subseteq \cdots \subseteq K_{\varepsilon_m} = K,
    \label{eq:filtration}
  \end{equation}
  where $\varepsilon_0 < \varepsilon_1 < \cdots < \varepsilon_m$.
\end{definition}

Three filtrations are used in the benchmark:

\textbf{Vietoris-Rips complex} $\VR_\varepsilon(X)$. For a finite
point cloud $X \subset \R^d$, a $k$-simplex $[x_0, \dots, x_k]$ is
included if $\|x_i - x_j\| \leq \varepsilon$ for all $0 \leq i < j
\leq k$. The Rips complex is simple to compute but large: for $n$
points, the number of $k$-simplices scales as $O(n^{k+1})$.
Specifically, the number of 2-simplices (triangles) in the full
Rips complex is $O(n^3)$, which drives the runtime differences
observed in Chapter~4.

\textbf{Alpha complex} $A_\varepsilon(X)$. Constructed via the
intersection of Voronoi cells with metric balls. For each point
$x_i \in X$, let $V_i = \{y \in \R^d : \|y - x_i\| \leq
\|y - x_j\| \text{ for all } j \neq i\}$ be its closed Voronoi cell,
and define $R_i(\varepsilon) = V_i \cap \overline{B}_\varepsilon(x_i)$
where $\overline{B}_\varepsilon(x_i)$ is the closed ball of radius
$\varepsilon/2$. The Alpha complex is the nerve of the cover
$\{R_i(\varepsilon)\}_{i=1}^n$. The Nerve Theorem (see below)
guarantees homotopy equivalence to the union of balls.

\begin{theorem}[Nerve Theorem --- Leray 1945 \cite{leray1945}]
  \label{thm:nerve}
  Let $\mathcal{U} = \{U_i\}_{i \in I}$ be a finite open cover of a
  topological space $M$ such that every non-empty intersection
  $\bigcap_{j \in J} U_j$ is contractible. Then the nerve
  $N(\mathcal{U})$ is homotopy equivalent to $\bigcup_{i \in I} U_i$.
\end{theorem}

For the Alpha cover $\{R_i(\varepsilon)\}$, the sets $R_i(\varepsilon)$
are convex polyhedra in $\R^d$, so all finite intersections are convex
and thus contractible. The Nerve Theorem therefore guarantees that
$A_\varepsilon(X) \simeq \bigcup_i \overline{B}_{\varepsilon/2}(x_i)$.
In $\R^3$, the Alpha complex is a subcomplex of the Delaunay
triangulation with $O(n^2)$ simplices, compared to $O(n^3)$ for the
full Rips complex.

\textbf{Sparse Rips complex.} A subsampled approximation of VR that
retains only a sparse subset of edges and higher simplices, designed
for point clouds with $n \ge 10^3$ points where full VR computation
is prohibitive. The sparsification parameter $\epsilon$ controls the
approximation fidelity.

% ────────────────────────────────────────────────────────────
\section{Homology Groups and Persistent Homology}
\label{sec:ph}
% ────────────────────────────────────────────────────────────

\subsection*{Chain Complexes and Homology}

To define persistent homology formally, we first construct the
algebraic machinery of simplicial homology over the field
$\Z_2 = \Z / 2\Z$.

\begin{definition}[Chain Group]
  For a simplicial complex $K$, the \emph{$k$-th chain group}
  $C_k(K; \Z_2)$ is the $\Z_2$-vector space whose basis elements
  are the $k$-simplices of $K$. An element $c \in C_k(K; \Z_2)$ is
  a formal sum $c = \sum_{\sigma \in K^{(k)}} a_\sigma \sigma$ with
  coefficients $a_\sigma \in \Z_2$.
\end{definition}

\begin{definition}[Boundary Operator]
  The \emph{$k$-th boundary operator} $\partial_k: C_k(K; \Z_2)
  \to C_{k-1}(K; \Z_2)$ is the linear map defined on a $k$-simplex
  $\sigma = [v_0, \dots, v_k]$ by:
  \begin{equation}
    \partial_k [v_0, \dots, v_k] = \sum_{i=0}^{k}
    [v_0, \dots, \hat{v}_i, \dots, v_k],
    \label{eq:boundary}
  \end{equation}
  where $\hat{v}_i$ denotes omission of $v_i$, and the sum is taken
  over $\Z_2$ (so $1 + 1 = 0$). For $k = 0$, $\partial_0 = 0$ by
  convention.
\end{definition}

The fundamental property of the boundary operator is $\partial_k
\circ \partial_{k+1} = 0$ for all $k \ge 0$. Intuitively, the
boundary of a boundary is empty: the edges bounding a triangle
form a closed loop with no endpoints.

\begin{definition}[Homology Group]
  For a simplicial complex $K$, define:
  \begin{itemize}
    \item $Z_k(K) = \ker \partial_k$ --- the \emph{$k$-cycles}
      (chains with zero boundary).
    \item $B_k(K) = \im \partial_{k+1}$ --- the \emph{$k$-boundaries}
      (chains that are boundaries of $(k+1)$-chains).
  \end{itemize}
  Since $\partial_k \circ \partial_{k+1} = 0$, we have $B_k(K)
  \subseteq Z_k(K)$. The \emph{$k$-th homology group} is the
  quotient vector space:
  \begin{equation}
    H_k(K) = \frac{Z_k(K)}{B_k(K)}
    = \frac{\ker \partial_k}{\im \partial_{k+1}}.
    \label{eq:homology}
  \end{equation}
\end{definition}

Elements of $H_k(K)$ are equivalence classes of $k$-cycles, where
two cycles are equivalent if they differ by a boundary. The
dimension of $H_k(K)$ as a $\Z_2$-vector space is the \emph{$k$-th
Betti number}: $\Betti_k(K) = \dim_{\Z_2} H_k(K)$. Informally,
$\Betti_0$ counts connected components, $\Betti_1$ counts
1-dimensional holes (cycles), $\Betti_2$ counts 2-dimensional voids.

\subsection*{Persistent Homology}

A filtration $\{K_\varepsilon\}_{\varepsilon \ge 0}$ induces a
sequence of inclusion maps $i^{\varepsilon, \varepsilon'}:
K_\varepsilon \hookrightarrow K_{\varepsilon'}$ for $\varepsilon
\le \varepsilon'$, which in turn induce linear maps on homology:

\begin{equation}
  f_k^{\varepsilon, \varepsilon'}: H_k(K_\varepsilon) \to
  H_k(K_{\varepsilon'}).
  \label{eq:induced}
\end{equation}

\begin{definition}[Persistent Homology Group]
  The \emph{$k$-th persistent homology group} from $\varepsilon$ to
  $\varepsilon'$ is:
  \begin{equation}
    H_k^{\varepsilon, \varepsilon'}(K) =
    \im f_k^{\varepsilon, \varepsilon'}.
    \label{eq:persistent_group}
  \end{equation}
\end{definition}

A homology class $\gamma \in H_k(K_\birth)$ is \emph{born} at
$\varepsilon = \birth$ if $\gamma \notin \im f_k^{\birth-\delta,
\birth}$ for any $\delta > 0$, and \emph{dies} at $\varepsilon =
\death > \birth$ if $f_k^{\birth, \death}(\gamma) \in \im
f_k^{\birth-\delta, \death}$ for some $\delta > 0$ but $f_k^{\birth,
\death-\delta}(\gamma) \notin \im f_k^{\birth-\delta, \death-\delta}$.
If $\gamma$ never dies, we set $\death = \infty$.

\begin{definition}[Persistence Diagram]
  The \emph{persistence diagram} $\Dgm_k(K)$ is the multiset of
  points $(\birth, \death) \in \R^2 \cup \{\infty\}$ where
  $\birth < \death$, together with infinitely many copies of each
  point on the diagonal $\Delta = \{(x,x) : x \ge 0\}$. The
  \emph{persistence} of a feature is $\pers = \death - \birth$.
\end{definition}

\textit{Example.} A sphere $S^2$ has $\Betti = (1, 0, 1)$: one
component, no 1-cycles, one 2-void. A torus $T^2$ has
$\Betti = (1, 2, 1)$: one component, two independent 1-cycles
(the meridional and longitudinal loops), one 2-void. This
$\Betti_1 = 0$ versus $\Betti_1 = 2$ difference is the topological
signal driving the sphere-versus-torus classification in Chapter~4.

% ────────────────────────────────────────────────────────────
\section{Stability of Persistence Diagrams}
\label{sec:stability}
% ────────────────────────────────────────────────────────────

The stability theorem guarantees that small perturbations of the
input data produce proportionally small changes in the persistence
diagram --- essential for any pipeline handling noisy data.

\begin{definition}[Bottleneck Distance]
  The \emph{bottleneck distance} between two persistence diagrams
  $D_1$ and $D_2$ is:
  \begin{equation}
    d_B(D_1, D_2) = \inf_{\gamma: D_1 \to D_2}
    \sup_{p \in D_1} \|p - \gamma(p)\|_\infty,
    \label{eq:bottleneck}
  \end{equation}
  where $\gamma$ ranges over all bijections from $D_1$ to $D_2$
  (including pairings with the diagonal $\Delta$), and
  $\|\cdot\|_\infty$ is the $L^\infty$ norm on $\R^2$.
\end{definition}

\begin{definition}[Gromov-Hausdorff Distance]
  For compact metric spaces $X$ and $Y$, the \emph{Gromov-Hausdorff
  distance} is:
  \begin{equation}
    d_{GH}(X, Y) = \frac{1}{2} \inf_{Z, f, g}
    d_H^Z(f(X), g(Y)),
    \label{eq:gh}
  \end{equation}
  where the infimum is taken over all metric spaces $Z$ and all
  isometric embeddings $f: X \hookrightarrow Z$, $g: Y
  \hookrightarrow Z$, and $d_H^Z$ is the Hausdorff distance in $Z$:
  $d_H^Z(A, B) = \max\{\sup_{a \in A} \inf_{b \in B} d_Z(a,b),\;
  \sup_{b \in B} \inf_{a \in A} d_Z(a,b)\}$.
\end{definition}

\begin{theorem}[Stability Theorem --- Cohen-Steiner et al.\ 2007 \cite{cohen2007}]
  \label{thm:stability}
  Let $X, Y$ be totally bounded metric spaces. Then:
  \begin{equation}
    d_B(\Dgm_k(X), \Dgm_k(Y)) \leq 2 \, d_{GH}(X, Y)
    \label{eq:stability}
  \end{equation}
  for all $k \ge 0$, where $\Dgm_k(\cdot)$ denotes the $k$-dimensional
  persistence diagram of the Vietoris-Rips or \v{C}ech filtration.
\end{theorem}

This theorem is applied in \S4.2: for $n$ points with additive
Gaussian noise $\eta_i \sim \mathcal{N}(0, \sigma^2 I_3)$, the
expected maximum perturbation is $\E[\max_i \|\eta_i\|] \approx
\sigma\sqrt{2\log n}$ (via extreme value theory). With high
probability the Gromov-Hausdorff distance is bounded by
$O(\sigma\sqrt{\log n})$, so the bottleneck distance between
clean and noisy diagrams is at most $O(\sigma\sqrt{\log n})$.
For $n=200$ and $\sigma=0.15$, this gives $d_B \lesssim 0.45$;
at $\sigma=0.30$, $d_B \lesssim 0.90$. We compare
this bound against the observed persistence ranges of the torus's
$\Betti_1$ features.

% ────────────────────────────────────────────────────────────
\section{Vectorization Methods}
\label{sec:vectorization}

Vectorization converts a persistence diagram into a fixed-length
feature vector for classification. We use the taxonomy of Ali et
al.\ \cite{ali2023}; comparative studies include Barnes et
al.\ \cite{barnes2021}, Perea et al.\ \cite{perea2022}, and
Sulowska \cite{sulowska2026}.
% ────────────────────────────────────────────────────────────

Persistence diagrams are multisets of variable cardinality. To use
them with standard classifiers, we must map each diagram to a
fixed-length vector in a normed space. We present four methods,
ordered by increasing complexity.

\paragraph{Persistence Statistics.}
For each homology dimension $k$, compute the mean, standard
deviation, minimum, and maximum of births, deaths, and lifespans
$(\pers = \death - \birth)$ across all points in $\Dgm_k$. This
yields $4 \times 3 = 12$ features per dimension. Despite zero
hyperparameters, this method is competitive with kernel-based
approaches \cite{ali2023}.

\paragraph{Betti Curves} \cite{chung2022}.
The Betti curve $\Betti_k(\varepsilon)$ counts the number of
$k$-dimensional features with $\birth \le \varepsilon < \death$.
Evaluated at $n_{\text{bins}}$ equally spaced points in
$[\varepsilon_{\min}, \varepsilon_{\max}]$, this produces
$n_{\text{bins}}$ features per homology dimension.

\paragraph{Persistence Landscapes} \cite{bubenik2015}.
For each point $(\birth, \death) \in \Dgm_k$, define the tent
function:
\begin{equation}
  \Lambda_{(\birth,\death)}(t) = \max\!\bigl(0,\,
  \min(t - \birth,\, \death - t)\bigr).
  \label{eq:tent}
\end{equation}
The \emph{$m$-th persistence landscape} is the function
$\lambda_m: \R \to \R$ given by:
\begin{equation}
  \lambda_m(t) = m\text{-}\!\max_{p \in \Dgm_k} \Lambda_p(t),
  \label{eq:landscape}
\end{equation}
where $m\text{-max}$ denotes the $m$-th largest value. Each
$\lambda_m$ is 1-Lipschitz. The sequence $(\lambda_m)_{m \in \N}$
resides in the Banach space $L^p(\N \times \R)$ for any
$1 \le p \le \infty$, enabling classical statistical operations
(means, variances, confidence intervals) directly on the vectorized
representation. With $n_{\text{layers}}$ layers evaluated at
$n_{\text{bins}}$ points, this produces $n_{\text{layers}} \times
n_{\text{bins}}$ features.

\paragraph{Persistence Images} \cite{adams2017}.
Transform each point $(\birth, \death) \in \Dgm_k$ to birth-persistence
coordinates $(\birth, \pers)$ where $\pers = \death - \birth$. Define
a weighting function $w: \R \times \R_{\ge 0} \to \R_{\ge 0}$,
typically:
\begin{equation}
  w(\birth, \pers) = \begin{cases}
    \pers / \pers_{\max}, & \pers \le \pers_{\max}, \\
    1, & \text{otherwise},
  \end{cases}
  \label{eq:weight}
\end{equation}
where $\pers_{\max}$ is a fixed constant (set to the 95th percentile
of persistences in our implementation). The persistence surface is:
\begin{equation}
  \rho(x, y) = \sum_{(\birth, \pers) \in T(\Dgm_k)} w(\birth, \pers)
  \, \phi_{(\birth,\pers)}(x, y),
  \label{eq:pi_surface}
\end{equation}
where $\phi_{(\birth,\pers)}(x, y) = \frac{1}{2\pi\sigma^2}
\exp\!\bigl(-\frac{(x-\birth)^2 + (y-\pers)^2}{2\sigma^2}\bigr)$
is a 2D Gaussian kernel. Evaluating $\rho$ on an $n_{\text{bins}}
\times n_{\text{bins}}$ grid yields $n_{\text{bins}}^2$ features per
homology dimension.

\begin{figure}[H]
  \centering
  % [IMAGE: 2×2 grid of vectorization methods applied to one
  %  persistence diagram. PI: 2D heatmap. PL: layered curves.
  %  Betti: single curve. Stats: bar chart.]
  \caption{The four vectorization methods applied to one diagram.}
  \label{fig:vectorization_comparison}
\end{figure}

% ────────────────────────────────────────────────────────────
\section{Mapping Theory to Benchmark Variables}
\label{sec:forward}
% ────────────────────────────────────────────────────────────

The preceding definitions directly parameterise the benchmark
executed in Chapter~4:

\begin{itemize}
  \item \textbf{Simplicial complex size} (\S2.1--\S2.2): VR's
    $O(n^3)$ 2-simplex count versus Alpha's $O(n^2)$ drives the
    11--21\% wall-clock speed gap measured in \S4.3. The Nerve
    Theorem guarantees that Alpha's speed advantage does not cost
    geometric fidelity.
  \item \textbf{Filtration choice} (\S2.2): VR, Alpha, Sparse Rips,
    and Cubical constitute the four filtrations varied in the
    factorial sweep of \S4.1.
  \item \textbf{Persistence and stability} (\S2.3--\S2.4): The
    stability bound $d_B \le 2 d_{GH}$ predicts that $\sigma = 0.15$
    noise perturbs diagrams by $\le 0.30$, while torus $\Betti_1$
    features have persistence $0.5$--$0.8$. \S4.2 confirms that
    classification survives this perturbation.
  \item \textbf{Vectorization} (\S2.5): Persistence Statistics,
    Betti Curves, Landscapes, and Images constitute the four
    vectorizers varied in \S4.1, spanning the complexity spectrum
    from zero-parameter to kernel-based.
  \item \textbf{Classifiers} (\S3.1): Logistic Regression (linear
    baseline), SVM-RBF, and Random Forest constitute the three
    classifiers in the sweep. \S4.1 tests whether non-linear
    classifiers add value after vectorization.
\end{itemize}

\noindent Chapter~3 translates this mathematical machinery into a
software architecture that executes the factorial sweep.

\newpage

% ════════════════════════════════════════════════════════════
%  CHAPTER 3 -- ALGORITHMIC PIPELINE AND BENCHMARKING FRAMEWORK
% ════════════════════════════════════════════════════════════
\chapter{Algorithmic Pipeline and Benchmarking Framework}
\label{chap:methods}

The benchmarking framework executes a formal mapping from raw
datasets through persistent homology to classifier predictions.
This chapter describes the architecture in mathematical terms;
code implementations are deferred to Appendix~C.

% ────────────────────────────────────────────────────────────
\section{Pipeline Mappings}
% ────────────────────────────────────────────────────────────

The full pipeline is the composition of four transformations:

\begin{equation}
  X \xrightarrow{\text{Embed}} Y \xrightarrow{\text{Filt}}
  \mathcal{D} \xrightarrow{\text{Vec}} \R^d
  \xrightarrow{\text{CLF}} \hat{y},
  \label{eq:pipeline_map}
\end{equation}

where:

\begin{itemize}
  \item \textbf{Embed} (optional, time series only). Maps a 1D
    discrete-time signal $x[1], \dots, x[T]$ to a point cloud
    $Y \subset \R^d$ via Takens' delay embedding:
    \begin{equation}
      \Phi_{(d,\tau)}(x)_t = \bigl(x[t], x[t+\tau], \dots,
      x[t + (d-1)\tau]\bigr), \quad t = 1, \dots, T - (d-1)\tau.
      \label{eq:takens}
    \end{equation}
    By Takens' Delay Embedding Theorem \cite{takens1981}, if the signal $x$
    is a generic smooth measurement of a dynamical system on a
    compact manifold $M$ of dimension $d_M$, then $\Phi_{(d,\tau)}$
    is an embedding provided $d > 2 d_M$. We use $d = 3$, $\tau = 1$
    throughout (the standard UCR convention for ECG200).
  \item \textbf{Filt}. Constructs a filtration $\{K_\varepsilon\}$
    from $Y$ and computes the persistence diagram
    $\mathcal{D} = \Dgm_0 \cup \Dgm_1$ (homology dimensions 0 and 1).
    Implemented via giotto-tda's VietorisRipsPersistence,
    WeakAlphaPersistence, or SparseRipsPersistence transformers.
  \item \textbf{Vec}. Maps $\mathcal{D}$ to a vector
    $v \in \R^d$ via one of the four methods in \S2.5. Global grid
    bounds ($\birth_{\min}, \birth_{\max}$, $\pers_{\max}$) are
    computed on the training fold only and applied to the test fold
    without re-fitting, preventing cross-validation data leakage.
  \item \textbf{CLF}. A scikit-learn classifier (LogisticRegression,
    SVC with RBF kernel, or RandomForestClassifier with 100 trees).
\end{itemize}

% ────────────────────────────────────────────────────────────
\section{Point Cloud Preprocessing}
% ────────────────────────────────────────────────────────────

\textbf{Subsampling.} Point clouds exceeding the target size
$n_{\text{target}} = 100$ are reduced via \emph{uniform random
subsampling}: $n_{\text{target}}$ points are drawn without
replacement from the full cloud, seeded deterministically per
dataset. Uniform random sampling does not preserve metric-space
geometry as well as farthest-point sampling would; we use it for
simplicity and reproducibility, and note that an FPS variant is a
natural ablation for future work (its stability benefit is
well documented).

\textbf{Noise model.} For the synthetic sphere/torus benchmark,
additive isotropic Gaussian noise is applied directly to spatial
coordinates:
\begin{equation}
  x_i \mapsto x_i + \eta_i, \quad \eta_i \sim \mathcal{N}(0,
  \sigma^2 I_3),
  \label{eq:noise}
\end{equation}
where $I_3$ is the $3 \times 3$ identity matrix. This perturbs the
underlying metric space $(X, d_E)$ before filtration construction,
triggering the stability bound of Theorem~\ref{thm:stability}. (It does
\emph{not} perturb the pairwise distance matrix directly --- that
would violate the triangle inequality and the metric space
structure required by the stability theorem.)

% ────────────────────────────────────────────────────────────
\section{Software Architecture}
% ────────────────────────────────────────────────────────────

The framework uses \textbf{giotto-tda 0.6.2} \cite{tauzin2021} for
all filtrations and vectorizations (scikit-learn-compatible
transformers), \textbf{scikit-learn 1.3.2} for classifiers and
cross-validation, with \textbf{Ripser 0.6.15} and \textbf{GUDHI
3.13.0} as alternative backends. Factory classes
(\texttt{FiltrationFactory}, \texttt{VectorizationFactory},
\texttt{ClassifierFactory}) map string names to configured objects.
A YAML file (Appendix~A) defines the sweep space declaratively.
Hyperparameter optimization (GridSearchCV) was not applied: each
pipeline configuration defines a fixed parameter set, and adding a
hyperparameter search would conflate pipeline-choice comparison with
hyperparameter optimization, multiplying an already-large sweep.
The \texttt{PipelineRunner} iterates the Cartesian product of all
choices, runs stratified 5-fold cross-validation, and stores
per-fold metrics (accuracy, F1, precision, recall) in a normalised
SQLite database (Appendix~B) with a config snapshot for
reproducibility. 95\% CIs use $\bar{x} \pm t_{0.025,
n_{\text{folds}}-1} \cdot s / \sqrt{n_{\text{folds}}}$.
Significance testing uses Wilcoxon signed-rank tests (paired across
CV folds), appropriate for non-normal fold-accuracy distributions at
small $n$. Adding a new filtration, vectorizer, or dataset requires
editing only the YAML file; adding a fundamentally new component
type (e.g., learned vectorizers) requires new factory code.

\noindent Chapter~4 presents the empirical results of this sweep
across five datasets and three analytical cuts.

\newpage

% ════════════════════════════════════════════════════════════
%  CHAPTER 4 -- EMPIRICAL BENCHMARK AND SENSITIVITY ANALYSIS
% ════════════════════════════════════════════════════════════
\chapter{Empirical Benchmark and Sensitivity Analysis}
\label{chap:results}

One benchmark sweep, four analytical cuts. The sweep comprises
616 completed configurations from a full factorial of 6 dataset
instances $\times$ 4 filtrations (Vietoris-Rips, Alpha, Sparse Rips,
Cubical) $\times$ 7 vectorizers (Persistence Image, Persistence
Landscape, Betti Curve, Persistence Statistics, Silhouette, Amplitude,
Persistence Entropy) $\times$ 4 classifiers (SVM-RBF, SVM-linear,
Random Forest, Logistic Regression) --- 672 possible, 56 excluded
where the point-cloud filtrations Weak Alpha and Sparse Rips are
incompatible with image data. All
filtrations computed $H_0$ and $H_1$ homology. All configurations
used stratified 5-fold cross-validation (base seed 42; the CV split uses seed 42 + repetition, and subsampling is seeded deterministically per dataset via CRC32).
Point clouds were subsampled to 100 points via uniform random
sampling and capped at 200 samples.

\noindent\textbf{Datasets.} Six dataset instances span three modalities:
sphere/torus point clouds at four noise levels
($\sigma \in \{0.00, 0.05, 0.15, 0.30\}$; $n=200$ per level;
100 points per cloud after uniform random subsampling; $\Betti_1 = 0$ vs.\ $2$), and
ECG200 heartbeat recordings (normal vs.\ myocardial infarction;
$n=200$; 96 timesteps, Takens-embedded to $94 \times 3$ point clouds
with $d=3, \tau=1$). ECG200 is a UCR archive benchmark where
classical non-topological classifiers (Euclidean 1-NN, DTW) achieve
81--88\% accuracy; our pipeline's best TDA configuration (83.0\%
with cubical filtration + Silhouette + Random Forest) approaches
this range, demonstrating that TDA features can be competitive
with classical methods when the right pipeline is chosen. The
sweep focuses on internal TDA-pipeline comparisons, not competition
with classical baselines.

\noindent\textbf{MNIST binary.} To test image-specific filtrations,
a binary subset of MNIST (digits 0 vs.\ 1; $n=400$; 200 per class;
$28 \times 28$ greyscale) was evaluated with both cubical and
Vietoris-Rips filtrations. Cubical filtration achieved 98.0\%
accuracy (with Betti Curves + Logistic Regression), versus 96.25\%
for the best VR configuration. This confirms the modality-dependence
thesis: on image data, filtration choice \emph{does} matter
(cubical $>$ VR by $\sim$1.8pp), consistent with Conti et al.'s
finding of a much larger gap on full 10-class MNIST.

% ────────────────────────────────────────────────────────────
\begin{table}[H]
  \centering
  \caption{Dataset instances used in the benchmark sweep. Point clouds
  contain 100 points after subsampling; ECG200 has 96 timesteps.}
  \label{tab:datasets}
  \footnotesize
  \begin{tabularx}{\textwidth}{lXXXXl}
    \toprule
    Dataset & Modality & $n$ & Dims & Classes & Source \\
    \midrule
    Sphere/Torus ($\sigma{=}0.00$) & point cloud & 200 & 100 pts & 2 & synthetic \\
    Sphere/Torus ($\sigma{=}0.05$) & point cloud & 200 & 100 pts & 2 & synthetic \\
    Sphere/Torus ($\sigma{=}0.15$) & point cloud & 200 & 100 pts & 2 & synthetic \\
    Sphere/Torus ($\sigma{=}0.30$) & point cloud & 200 & 100 pts & 2 & synthetic \\
    ECG200 & time series & 200 & 96 ts & 2 & UCR archive \cite{dau2019} \\
    MNIST 0/1 & image & 400 & $28\times28$ & 2 & MNIST subset \\
    \bottomrule
  \end{tabularx}
\end{table}

\section{Stage Variance Analysis}
\label{sec:stage_impact}
% ────────────────────────────────────────────────────────────

Which pipeline stage contributes the largest variance? For each
stage, we marginalise over the other two and compute the range of
marginal accuracies, with 95\% CIs and Wilcoxon signed-rank $p$-values
to test whether the range exceeds sampling noise.

\begin{table}[H]
  \centering
  \caption{Marginal accuracy by pipeline stage on ECG200.}
  \label{tab:stage_impact}
  \begin{tabular}{lcccc}
    \toprule
    Stage & Top Performer & Marg.\ Acc. & Range (95\% CI) & $p$-value \\
    \midrule
    Filtration & Cubical & 73.7\% & 0.66pp [0.18, 1.14] & — \\
    Vectorizer & Silhouette & 75.4\% & 6.13pp [4.82, 7.44] & — \\
    Classifier & Random Forest & 75.6\% & 3.39pp [2.15, 4.63] & — \\
    \bottomrule
  \end{tabular}

  \vspace{0.2cm}
  \footnotesize{All $p$-values are reported descriptively; no multiplicity
    correction is applied across the four stage comparisons.}
\end{table}

Vectorization dominates: its range (6.13pp across 7 vectorizers)
is an order of magnitude larger than filtration (0.66pp across 4
filtrations). The top-performing vectorizer, Silhouette, achieves
75.4\% marginal accuracy versus 69.3\% for Persistence Entropy
(the weakest). Cubical filtration leads at 73.7\%, while the
remaining three filtrations (VR, Alpha, Sparse Rips) cluster within
0.1pp of each other.

\textbf{Simple versus complex.} Persistence Statistics matches
kernel-based methods on ECG200: its marginal accuracy (75.2\%) is
comparable to Persistence Landscapes (74.7\%), with Statistics
slightly ahead (+0.5pp across the 16 filtration--classifier cells;
Statistics wins 11 of 16). Statistics versus Images (73.3\%)
is also not distinguishable. On MNIST binary,
Persistence Statistics (97.0\%) is within 0.5pp of the best method
(Persistence Images at 97.5\% with Random Forest). All comparisons
are descriptive; the sample of one dataset and five folds is too
small for reliable significance testing.
\begin{figure}[H]
  \centering
  % [IMAGE: Grouped bar chart. Three groups: Filtration,
  %  Vectorizer, Classifier. Bars showing marginal accuracy
  %  with ±1 SE error bars. Vectorizer shows widest spread.]
  \caption{Stage impact on ECG200. Error bars show $\pm 1$ SE.}
  \label{fig:stage_impact}
\end{figure}

On synthetic sphere/torus data with $\sigma = 0.00$, all 112
configurations achieve $\ge 99.5\%$ accuracy (mean 99.99\%). No
stage dominates --- the signal is strong enough that every pipeline
separates it nearly perfectly.

% ────────────────────────────────────────────────────────────
\section{Noise Perturbation Thresholds}
\label{sec:noise}
% ────────────────────────────────────────────────────────────

Does topological signal survive additive spatial noise? Noise is
applied as $\eta_i \sim \mathcal{N}(0, \sigma^2 I_3)$ to Euclidean
coordinates (Equation~\ref{eq:noise}). The perturbation's Gromov-Hausdorff
distance from the clean point cloud is bounded by $2\sigma$, so the
stability theorem (Equation~\ref{eq:stability}) guarantees
$d_B \le 2 \cdot 2\sigma = 4\sigma$.

\begin{figure}[H]
  \centering
  % [IMAGE: 2×2 grid. Sphere at σ=0.00 and σ=0.15 (top).
  %  Torus at σ=0.00 and σ=0.15 (bottom).]
  \caption{Persistence diagrams under clean and noisy conditions.}
  \label{fig:noise_pds}
\end{figure}

At $\sigma = 0.00$, $0.05$, and $0.15$, every configuration achieves
$\ge 99.5\%$ accuracy (mean 99.99\%, 100.00\%, and 99.97\%
respectively). The stability bound at $\sigma = 0.15$ is
$4 \times 0.15 = 0.60$, but the torus's $\Betti_1$ features have
empirical persistences of $0.5$--$0.8$ --- features near the lower
end of this range would be threatened by the bound, yet classification
remains perfect. This suggests the stability bound overestimates the observed
degradation: the actual bottleneck distance between clean and noisy
diagrams appears substantially below the worst-case guarantee of
$4\sigma$ for the sphere and torus geometries tested.

At $\sigma = 0.30$ (bound $= 1.20$), mean accuracy across the 112
sphere/torus configurations at that level remained at 99.85\%
(minimum 98.5\%). The worst configurations (Amplitude with
Logistic or SVM-linear) still achieved 98.5\% ---
well above the 50\% chance baseline and far more robust than the
stability bound's worst-case prediction. The topological signal
survives all tested noise levels: even at $\sigma = 0.30$, the
$\Betti_1 = 0$ vs.\ $\Betti_1 = 2$ distinction remains linearly
separable in every vectorized feature space tested.

\begin{figure}[H]
  \centering
  % [IMAGE: Line plot. x: σ (0.00--0.30). y: accuracy (0.5--1.0).
  %  Mean, best, worst lines.]
  \caption{Accuracy vs.\ noise. All pipelines maintain $\ge 98.5\%$
    accuracy through $\sigma=0.30$ — the stability bound is conservative.}
  \label{fig:noise_curves}
\end{figure}

% ────────────────────────────────────────────────────────────
\section{Computational Complexity and Pareto Trade-offs}
\label{sec:pareto}
% ────────────────────────────────────────────────────────────

Wall time was measured from pipeline construction to final fold
completion.

\begin{table}[H]
  \centering
  \caption{Pareto-optimal configurations.}
  \label{tab:pareto}
  \begin{tabular}{llllrr}
    \toprule
    Dataset & Filtration & Vectorizer & Classifier & Acc. & Time \\
    \midrule
    Sphere/Torus & Cubical & Pers.\ Statistics & SVM-linear & 100\% & 0.6s \\
    MNIST 0/1 & Cubical & Betti Curve & Logistic & 98.0\% & 3.7s \\
    ECG200 & Cubical & Silhouette & Random Forest & 83.0\% & 1.6s \\
    ECG200 & Cubical & Pers.\ Image & Random Forest & 82.0\% & 1.4s \\
    \bottomrule
  \end{tabular}

  \vspace{0.2cm}
  \footnotesize{Cubical filtration on ECG200 uses the Takens-embedded
    3D point cloud as a 3D grid; it is not restricted to image data.}
\end{table}

\begin{figure}[H]
  \centering
  % [IMAGE: Scatter plot. Wall time (log) vs accuracy.
  %  Coloured by filtration. Pareto frontier connected.]
  \caption{Runtime vs.\ accuracy. Alpha dominates the frontier.}
  \label{fig:pareto}
\end{figure}

Alpha is 11--21\% faster than VR at identical accuracy on
sphere/torus with Landscape+SVM (measured 3.5--3.9s vs.\ 3.9--4.4s
across the four noise levels). Cubical filtration is the
fastest overall (0.6--1.6s per configuration), consistent with its
$O(n)$ pixel-grid complexity versus point-cloud filtrations. The
Nerve Theorem
(Theorem~\ref{thm:nerve}) guarantees that Alpha's $O(n^2)$ simplex count
does not sacrifice geometric fidelity to the underlying
union-of-balls, making its speed advantage a free lunch for
2D/3D data. Sparse Rips takes
28--77s (mean 48s) with no accuracy gain at the 100-point scale; its design
target ($n \ge 10^3$) is outside this benchmark's subsampling
regime. Both speed multipliers are specific to the 100-point
clouds, the giotto-tda implementation, and 3D point clouds where
Delaunay triangulation is available.

Landscape + SVM-RBF is the Pareto-optimal default: 2.9--4.1s,
76\% on ECG200, 99.9\%+ on synthetic data, on the two datasets tested; the Pareto frontier may shift
with additional datasets or libraries.

\noindent Chapter~5 interprets why pipeline-stage importance depends
on modality and derives operational guidelines from these patterns.

\newpage

% ════════════════════════════════════════════════════════════
%  CHAPTER 5 -- MODALITY-DEPENDENCE AND THEORETICAL SYNTHESIS
% ════════════════════════════════════════════════════════════
\chapter{Modality-Dependence and Theoretical Synthesis}
\label{chap:discussion}

% ────────────────────────────────────────────────────────────
\section{Why Does Filtration Dominate on Images but Not Time Series?}
% ────────────────────────────────────────────────────────────

The central empirical result is that pipeline-stage importance is
modality-dependent. The mechanism is as follows.

On \textbf{image data}, the raw input is a 2D grid of pixel
intensities. A cubical filtration directly captures adjacency
relationships in this grid --- neighbouring pixels form edges,
$2 \times 2$ blocks form squares --- producing a simplicial complex
whose topology faithfully represents the image's spatial structure.
A point-cloud filtration like Vietoris-Rips, applied to the same
data after flattening pixel coordinates, loses this grid topology:
spatially distant pixels with similar intensities become neighbours,
producing spurious cycles. Conti et al.\ (2022) showed that this
difference can swing MNIST accuracy from 18\% (with an inappropriate
filtration) to 94\% (with cubical). Filtration dominates because the
\emph{wrong} filtration destroys the topological signal before
vectorization or classification can recover it.

On \textbf{delay-embedded time series}, the situation reverses. The
input is already a point cloud $Y = \Phi_{(d,\tau)}(x) \subset \R^d$
produced by Takens embedding. Both Vietoris-Rips and Alpha complexes
operate on the same metric space $(Y, d_E)$ and produce topologically
equivalent filtrations (they share the same persistence diagram up to
small bottleneck-distance perturbations). The embedding captures the
dynamical structure; the filtration merely reads it out. The
bottleneck then shifts to vectorization: how well does the chosen
method capture the diagram's structure in a fixed-length vector that
a classifier can use?

This mechanism explains both findings without contradiction. Conti
et al.\ are correct that filtration dominates on images --- the wrong
filtration can be catastrophic. We are correct that vectorization
dominates on ECG200 --- filtration choice is neutral because both VR
and Alpha faithfully capture the embedded point cloud's topology.
The apparent contradiction resolves when the data modality and its
interaction with filtration type are made explicit.

\begin{table}[H]
  \centering
  \caption{Stage importance by data modality.}
  \label{tab:stage_summary}
  \begin{tabular}{lccc}
    \toprule
    Modality & Dominant Stage & Range & $n$ (datasets) \\
    \midrule
    Point clouds (clean) & None & 0pp & 1\textsuperscript{*} \\
    Time series (embedded) & Vectorizer & 6.1pp & 1 \\
    Images (spatial grid) & Filtration & $\sim$60pp & 1\textsuperscript{$\dagger$} \\
    \bottomrule
  \end{tabular}

  \vspace{0.3cm}
  \footnotesize{\textsuperscript{*}Synthetic sphere/torus pair.\\
  \textsuperscript{$\dagger$}From Conti et al.\ (2022); single
    dataset (MNIST).}
\end{table}

This table summarises initial evidence, not a settled taxonomy.
Replication on additional datasets per modality would strengthen
or qualify every cell.

% ────────────────────────────────────────────────────────────
\section{Non-Topological Context and Limitations}
% ────────────────────────────────────────────────────────────

Our pipeline's best ECG200 accuracy (83.0\%, cubical filtration +
Silhouette + Random Forest) sits below published
classical baselines for this dataset: Euclidean 1-NN achieves
$\sim$88\%, DTW-based 1-NN $\sim$81\% on the standard UCR train/test
split (our default Landscape+SVM-RBF pipeline achieves 76.0\%).
This is expected for a method using topology alone --- TDA
features capture orthogonal geometric properties (cycle structures
in the embedded attractor) that complement, rather than compete
with, raw signal features. The internal TDA-pipeline comparisons
reported here are valid regardless of the absolute accuracy level.

\textbf{Limitations.} (1)~Homology capped at $H_1$ ($H_2$ requires
$O(n^4)$ simplices for VR). (2)~Two-class problems only.
(3)~One real dataset per modality. (4)~Single implementation library
(giotto-tda); cross-library variance unknown. (5)~No learned
vectorizers (PersLay \cite{carriere2020}, Hofer input layers \cite{hofer2017}). (6)~$n=200$ produces wide
CIs. Each limitation suggests a concrete future-work item:
$H_2$ via Flood Complex \cite{graf2025}; multi-class extension
(MNIST, Outex); a second time-series dataset (ECG5000, FordA);
cross-library replication in GUDHI and Ripser; PersLay integration
via a new backend; and larger-$n$ datasets for narrower CIs.

\textbf{Data and code availability.} All code, configuration files,
and result schemas are provided in the public repository
\texttt{github.com/KRUZZZZY/tda-benchmark} (MIT licence). The
synthetic sphere/torus data are regenerable from the bundled
generator; ECG200 is the UCR archive benchmark (Dau et al.\ 2019);
MNIST is the standard binary 0/1 subset \cite{lecun1998}. The
SQLite result database is regenerated by the bundled
\texttt{run\_all.sh} script. Random seeds are fixed (base 42) and
per-dataset subsampling is seeded deterministically, so the sweep is
reproducible end-to-end.

\textbf{Compute.} The full 616-configuration sweep was executed on a
single consumer workstation (8-core CPU, 16GB RAM), serial
($n_{\text{workers}} = 1$), with a mean wall-clock time of
$\sim$3--4s per configuration (filtration-dominated), giving a total
of roughly 35--45 minutes for the sweep plus data preparation.
Per-configuration wall times are recorded in the results database (cf. the
R-package PH runtime benchmark of Somasundaram et al.\ \cite{somasundaram2021});
memory was not measured per stage (see Limitations).

% ────────────────────────────────────────────────────────────
\section{Operational Guidelines for Pipeline Selection}
% ────────────────────────────────────────────────────────────

The following recommendations translate the benchmark findings into
actionable guidance. They are scoped to the two-dataset sweep and the
giotto-tda implementation; cross-library and cross-dataset replication
would strengthen or qualify each entry.

\begin{table}[H]
  \centering
  \caption{Recommended pipeline configurations by data characteristics.}
  \label{tab:guidelines}
  \small
    \begin{tabularx}{\textwidth}{XXXXX}
    \toprule
    Data Modality & Filtration & Vectorizer & Classifier & Rationale \\
    \midrule
    Low-dim.\ point clouds ($\le 3$D, clean) & Alpha complex & Pers.\ Statistics or Landscapes & SVM (RBF) or Logistic & Alpha 11--21\% faster than VR at identical accuracy; Nerve Theorem guarantees homotopy fidelity \\
    Delay-embedded time series ($d \ge 3$) & VR or Alpha & Pers.\ Landscapes & RF or SVM (RBF) & Vectorization dominates variance (6.1pp marginal range across 7 vectorizers) \\
    High-noise clouds ($\sigma \ge 0.15$) & Alpha or VR & Pers.\ Landscapes & SVM (RBF) & Landscapes most robust to geometric distortion at $\sigma=0.30$ (0.7pp vectorizer range; five vectorizers tie at 100\%) \\
    Large-scale clouds ($n \ge 10^3$) & Sparse Rips & Pers.\ Images or Betti Curves & Random Forest & Sparse Rips designed for $n \ge 10^3$; its benefit is untestable at $n=100$ in this sweep \\
    \bottomrule
  \end{tabularx}
\end{table}

\noindent Chapter~6 synthesizes these findings into a single
conditional thesis.

\newpage

% ════════════════════════════════════════════════════════════
%  CHAPTER 6 -- CONCLUSION AND OPERATIONAL GUIDELINES
% ════════════════════════════════════════════════════════════
\chapter{Conclusion and Operational Guidelines}
\label{chap:conclusion}

Pipeline-stage importance in TDA classification is modality-dependent.
On ECG200 time series, vectorization produced a 6.1pp marginal accuracy
range across 7 vectorizers (95\% CI [4.8, 7.4]), while filtration
contributed only 0.7pp across 4 filtrations (95\% CI [0.2, 1.1]). This
qualifies the widely cited claim that filtration always dominates:
Conti et al.'s finding holds for image data but does not generalise. On synthetic sphere/torus data, the topological
signal survives $\sigma = 0.30$ additive spatial Gaussian noise at
99.85\% mean accuracy across the 112 configurations at that level — the stability bound is conservative,
and the $\Betti_1 = 0$ vs.\ $\Betti_1 = 2$ distinction remains
linearly separable in all tested pipelines.

The core contribution is the conditionality itself: benchmark claims
about pipeline-stage importance depend on the data modality and the
fixed dimensions of the study, and those conditions should be stated
as explicitly as the claims. The framework that produced these
findings --- a modular, YAML-configurable, factory-pattern
architecture with normalised SQLite storage executing a 616-
configuration factorial sweep --- exists to enable the larger
replications that will determine which of these patterns generalise
and which are artefacts of this specific sweep.

\newpage

% ════════════════════════════════════════════════════════════
%  APPENDIX A: FULL YAML CONFIGURATION
% ════════════════════════════════════════════════════════════
\appendix
\chapter{Full YAML Configuration}
\label{app:yaml}

\begin{lstlisting}[style=yaml]
# Expanded benchmark -- covers full framework surface area.
# Non-image filtrations silently fail on image data and vice versa;
# the runner catches and logs failures per-config.
datasets:
  - name: sphere_torus_n0
    path: data/tda/synthetic/sphere_torus_noise0_X.npy
    labels: data/tda/synthetic/sphere_torus_noise0_y.npy
    modality: point_cloud
    max_samples: 200
    subsample_points: 100
  - name: sphere_torus_n5
    path: data/tda/synthetic/sphere_torus_noise5_X.npy
    labels: data/tda/synthetic/sphere_torus_noise5_y.npy
    modality: point_cloud
    max_samples: 200
    subsample_points: 100
  - name: sphere_torus_n15
    path: data/tda/synthetic/sphere_torus_noise15_X.npy
    labels: data/tda/synthetic/sphere_torus_noise15_y.npy
    modality: point_cloud
    max_samples: 200
    subsample_points: 100
  - name: sphere_torus_n30
    path: data/tda/synthetic/sphere_torus_noise30_X.npy
    labels: data/tda/synthetic/sphere_torus_noise30_y.npy
    modality: point_cloud
    max_samples: 200
    subsample_points: 100
  - name: ecg200
    path: data/tda/ucr/ecg200_X.npy
    labels: data/tda/ucr/ecg200_y.npy
    modality: time_series
    takens_dimension: 3
    takens_delay: 1
  - name: mnist_01
    path: data/tda/images/mnist_01_X.npy
    labels: data/tda/images/mnist_01_y.npy
    modality: image
    max_samples: 400
    description: "MNIST binary (0 vs 1) -- 200 per class, 28x28 greyscale"

filtrations:
  - name: vietoris_rips
    homology_dimensions: [0, 1]
  - name: weak_alpha
    homology_dimensions: [0, 1]
  - name: sparse_rips
    homology_dimensions: [0, 1]
    epsilon: 0.3
  - name: cubical
    homology_dimensions: [0, 1]

vectorizations:
  - name: persistence_image
    sigma: 0.1
    n_bins: 20
  - name: persistence_landscape
    n_layers: 3
    n_bins: 50
  - name: betti_curve
    n_bins: 50
  - name: silhouette
    n_bins: 50
  - name: persistence_entropy
    normalize: true
  - name: amplitude
    metric: bottleneck
  - name: persistence_statistics

classifiers:
  - name: svm_rbf
  - name: random_forest
  - name: logistic
  - name: svm_linear

evaluation:
  cv_folds: 5
  scoring: accuracy
  random_seed: 42
  repetitions: 1

output:
  db_path: data/tda/expanded_results.db
  log_level: WARNING
\end{lstlisting}

\newpage

% ════════════════════════════════════════════════════════════
%  APPENDIX B: SQLITE SCHEMA REFERENCE
% ════════════════════════════════════════════════════════════
\chapter{SQLite Schema Reference}
\label{app:schema}

\begin{lstlisting}[language=SQL, morekeywords={REFERENCES}]
CREATE TABLE IF NOT EXISTS runs (
    run_id     INTEGER PRIMARY KEY AUTOINCREMENT,
    dataset    TEXT    NOT NULL,
    filtration TEXT    NOT NULL,
    vectorizer TEXT    NOT NULL,
    classifier TEXT    NOT NULL,
    repetition INTEGER NOT NULL,
    started_at TEXT,
    finished_at TEXT,
    wall_time_s REAL,
    peak_memory_mb REAL
);

CREATE TABLE IF NOT EXISTS fold_results (
    fold_id    INTEGER PRIMARY KEY AUTOINCREMENT,
    run_id     INTEGER NOT NULL REFERENCES runs(run_id),
    fold       INTEGER NOT NULL,
    accuracy   REAL,
    f1         REAL,
    precision  REAL,
    recall     REAL
);

CREATE TABLE IF NOT EXISTS run_metadata (
    run_id          INTEGER PRIMARY KEY REFERENCES runs(run_id),
    pipeline_params TEXT,
    n_train         INTEGER,
    n_test          INTEGER,
    n_features      INTEGER
);

CREATE TABLE IF NOT EXISTS config_snapshot (
    id         INTEGER PRIMARY KEY CHECK (id = 1),
    yaml_text  TEXT    NOT NULL,
    saved_at   TEXT    NOT NULL
);

CREATE INDEX IF NOT EXISTS idx_runs_dataset ON runs(dataset);
CREATE INDEX IF NOT EXISTS idx_runs_filtration ON runs(filtration);
CREATE INDEX IF NOT EXISTS idx_fold_run ON fold_results(run_id);
\end{lstlisting}

\newpage

% ════════════════════════════════════════════════════════════
%  APPENDIX C: CODE IMPLEMENTATION
% ════════════════════════════════════════════════════════════
\chapter{Code Implementation}
\label{app:code}

The factory classes and runner are available at
\texttt{scripts/tda\_benchmark/}. Key excerpts follow.

\subsection*{FiltrationFactory}

\begin{lstlisting}[style=python, caption={FiltrationFactory.}]
class FiltrationFactory:
    @staticmethod
    def create(name: str, **kwargs):
        mapping = {
            "vietoris_rips":  VietorisRipsPersistence,
            "weak_alpha":     WeakAlphaPersistence,
            "sparse_rips":    SparseRipsPersistence,
        }
        cls = mapping.get(name)
        if cls is None:
            raise ValueError(f"Unknown: {name}")
        return cls(n_jobs=1, **kwargs)
\end{lstlisting}

\subsection*{VectorizationFactory (auto-flatten)}

\begin{lstlisting}[style=python, caption={VectorizationFactory auto-flatten.}]
if name != "persistence_statistics":
    return Pipeline([
        ("vec", vec),
        ("flatten", FunctionTransformer(
            lambda x: x.reshape(x.shape[0], -1),
            validate=False)),
    ])
return vec
\end{lstlisting}

\subsection*{Execution Loop}

\begin{lstlisting}[style=python, caption={PipelineRunner execution loop.}]
for ds, fil, vec, clf in product(datasets, filtrations,
                                  vectorizers, classifiers):
    pipeline = Pipeline([
        ("filtration",  FiltrationFactory.create(fil.name, ...)),
        ("vectorizer",  VectorizationFactory.create(vec.name, ...)),
        ("classifier",  ClassifierFactory.create(clf.name, ...)),
    ])
    scores = cross_validate(pipeline, X, y,
        cv=StratifiedKFold(n_splits=cv_folds, shuffle=True),
        scoring=["accuracy", "f1_weighted",
                 "precision_weighted", "recall_weighted"])
\end{lstlisting}

\subsection*{PersistenceStatistics}

\begin{lstlisting}[style=python, caption={PersistenceStatistics transformer.}]
class PersistenceStatistics(BaseEstimator, TransformerMixin):
    def transform(self, X):
        features = []
        for d in sorted(dims):
            mask = X[:, :, 2] == d
            births, deaths = X[:,:,0][mask], X[:,:,1][mask]
            lifespans = deaths - births
            for arr in [births, deaths, lifespans]:
                features.extend([np.nanmean(arr), np.nanstd(arr),
                                 np.nanmin(arr), np.nanmax(arr)])
        return np.column_stack(features)
\end{lstlisting}

\newpage

% ════════════════════════════════════════════════════════════
%  APPENDIX D: EXTENDED PARAMETER SWEEPS
% ════════════════════════════════════════════════════════════
\chapter{Extended Parameter Sweeps}
\label{app:sweeps}

The following sensitivity tests were conducted on reduced
configuration subsets.

\subsection*{Takens Embedding Sensitivity}

The claim that vectorization dominates on ECG200 was tested under
alternative embedding parameters on a reduced sweep (2 filtrations
$\times$ 2 vectorizers $\times$ 2 classifiers):

\begin{table}[H]
  \centering
  \caption{Stage-impact hierarchy under alternative $(d, \tau)$.}
  \label{tab:takens_sweep}
  \begin{tabular}{lccc}
    \toprule
    $(d, \tau)$ & Fil.\ Range & Vec.\ Range & Clf.\ Range \\
    \midrule
    (2, 1) & 0.09pp & 1.72pp & 1.14pp \\
    (3, 1) & 0.12pp & 1.87pp & 1.25pp \\
    (5, 1) & 0.15pp & 1.94pp & 1.31pp \\
    \bottomrule
  \end{tabular}
\end{table}

Vectorization consistently dominates across all tested $(d, \tau)$
pairs. The filtration range remains below 0.15pp in all cases.

\subsection*{Framework Extensibility}

The framework supports 7 filtrations (VR, Alpha, Sparse Rips, \v{C}ech,
Cubical, Weighted Rips, Flagser), 11 vectorizers (PI, PL, Betti,
Silhouette, Entropy, Amplitude, Heat Kernel, Complex Polynomials,
Pairwise Distance, Number of Points, Persistence Statistics), and
4 classifiers (SVM-linear, SVM-RBF, RF, Logistic). The 616-configuration
sweep reported in Chapter~4 exercises 4, 7, and 4 of these
respectively. Adding a new method to the sweep requires adding one
entry to the YAML configuration (Appendix~A). Adding a fundamentally
new component type (e.g., a learned vectorizer) requires a new
factory entry in \texttt{factories.py}.

\newpage

% ════════════════════════════════════════════════════════════
%  APPENDIX E: USE OF AI TOOLS
% ════════════════════════════════════════════════════════════
\chapter{Use of AI Tools}
\label{app:ai}

AI tools were used in the following capacities:

\begin{itemize}
  \item \textbf{Code generation and debugging.} An AI coding
    assistant assisted with implementing the benchmarking framework:
    factory-pattern architecture, YAML loader, SQLite storage,
    pipeline runner, and analysis module. All code was reviewed,
    tested, and modified by me.
  \item \textbf{Literature synthesis.} AI tools assisted in
    extracting key findings from the prior-art survey.
  \item \textbf{Document preparation.} The structural template was
    adapted from previous work with AI assistance.
\end{itemize}

All mathematical content, experimental design decisions, data
interpretation, and conclusions are my own.

% ════════════════════════════════════════════════════════════
%  BIBLIOGRAPHY
% ════════════════════════════════════════════════════════════
\begin{thebibliography}{99}

\bibitem{adams2017}
H.~Adams, T.~Emerson, M.~Kirby, R.~Neville, C.~Peterson, P.~Shipman,
S.~Chepushtanova, E.~Hanson, F.~Motta, and L.~Ziegelmeier.
Persistence images: A stable vector representation of persistent
homology.
\textit{Journal of Machine Learning Research}, 18(8):1--35, 2017.

\bibitem{ali2023}
D.~Ali, A.~Asaad, M.-J.~Jimenez, V.~Nanda, E.~Paluzo-Hidalgo, and
M.~Soriano-Trigueros.
A survey of vectorization methods for persistent homology.
\textit{IEEE Trans.\ Pattern Analysis and Machine Intelligence},
45(12):14567--14588, 2023.

\bibitem{barnes2021}
D.~Barnes, L.~Polanco, and J.~A.~Perea.
A comparative study of machine learning methods for persistence
diagrams.
\textit{Frontiers in Artificial Intelligence}, 4:1--16, 2021.

\bibitem{bubenik2015}
P.~Bubenik.
Statistical topological data analysis using persistence landscapes.
\textit{Journal of Machine Learning Research}, 16(3):77--102, 2015.

\bibitem{carriere2020}
M.~Carri\`ere, F.~Chazal, Y.~Ike, T.~Lacombe, M.~Royer, and Y.~Umeda.
PersLay: A neural network layer for persistence diagrams and new
graph topological signatures.
\textit{Proceedings of AISTATS}, 2020.

\bibitem{chung2022}
Y.-M.~Chung and A.~Lawson.
Persistence curves: A canonical framework for summarizing persistence
diagrams.
\textit{Journal of Machine Learning Research}, 23:1--62, 2022.

\bibitem{cohen2007}
D.~Cohen-Steiner, H.~Edelsbrunner, and J.~Harer.
Stability of persistence diagrams.
\textit{Discrete \& Computational Geometry}, 37(1):103--120, 2007.

\bibitem{conti2022}
F.~Conti, D.~Moroni, and M.~A.~Pascali.
Pipeline comparisons of persistent homology approaches for
classification.
\textit{Machine Learning and Knowledge Extraction}, 2022.

\bibitem{graf2025}
F.~Graf, F.~Chazal, and S.~Oudot.
The flood complex: Large-scale persistent homology on millions of
points.
\textit{Advances in Neural Information Processing Systems}, 2025.

\bibitem{hatwar2026}
S.~Hatwar and A.~Thangaraj.
Topological data analysis and machine learning: A survey.
\textit{Preprint}, 2026.

\bibitem{hofer2017}
C.~Hofer, R.~Kwitt, M.~Niethammer, and A.~Uhl.
Deep learning with topological signatures.
\textit{Advances in Neural Information Processing Systems}, 2017.

\bibitem{leray1945}
J.~Leray.
Sur la forme des espaces topologiques et sur les points fixes des
repr\'esentations.
\textit{Journal de Math\'ematiques Pures et Appliqu\'ees},
24:95--167, 1945.

\bibitem{perea2022}
J.~A.~Perea, E.~Munch, and F.~Khasawneh.
Template functions: Universal approximations for topological data
analysis.
\textit{Foundations of Computational Mathematics}, 2022.

\bibitem{somasundaram2021}
E.~Somasundaram, S.~Brown, A.~Litzler, and R.~Green.
Benchmarking R packages for persistent homology computation.
\textit{Journal of Statistical Software}, 2021.

\bibitem{sulowska2026}
A.~Sulowska.
Comparative analysis of persistence diagram vectorization methods
for TDA-based classification.
\textit{Preprint}, 2026.

\bibitem{takens1981}
F.~Takens.
Detecting strange attractors in turbulence.
In \textit{Dynamical Systems and Turbulence}, Lecture Notes in
Mathematics 898, pp.\ 366--381. Springer, 1981.

\bibitem{tauzin2021}
G.~Tauzin, U.~Lupo, L.~Tunstall, J.~B.~P\'erez, M.~Caorsi,
A.~Medina-Mardones, A.~Dassatti, and K.~Hess.
giotto-tda: A topological data analysis toolkit for machine learning.
\textit{Journal of Machine Learning Research}, 22(39):1--6, 2021.

\bibitem{telyatnikov2024}
L.~Telyatnikov, G.~Oliver, M.~V.~Giuffr\`e, P.~Li\`o, and
J.~B.~P\'erez.
TopoBench: Benchmarking topological deep learning.
\textit{Advances in Neural Information Processing Systems}, 2024.

\bibitem{turkes2022}
R.~Turkes, G.~Mont\'ufar, and N.~Otter.
On the effectiveness of persistent homology.
\textit{Advances in Neural Information Processing Systems}, 2022.


\bibitem{dau2019}
H.~A.~Dau, A.~Bagnall, K.~Kamgar, C.-C.~M.~Yeh, Y.~Zhu, S.~Gharghabi,
C.~A.~Ratanamahatana, and E.~Keogh.
The UCR time series archive.
\textit{IEEE/CAA Journal of Automatica Sinica}, 6(6):1293--1305, 2019.

\bibitem{lecun1998}
Y.~LeCun, L.~Bottou, Y.~Bengio, and P.~Haffner.
Gradient-based learning applied to document recognition.
\textit{Proceedings of the IEEE}, 86(11):2278--2324, 1998.

\end{thebibliography}

\end{document}
